\documentclass[11pt]{article}
\newcommand{\SheetCode}{Theory 03}
\usepackage[margin=25mm]{geometry}
\usepackage[T1]{fontenc}
\usepackage{lmodern}
\DeclareUnicodeCharacter{2605}{$\star$}
\usepackage{microtype}
\usepackage{amsmath,amssymb,mathtools,bm}
\usepackage{enumitem}
\usepackage{xcolor}
\usepackage{fancyhdr}
\usepackage{hyperref}
\usepackage[most]{tcolorbox}
\definecolor{agiBlue}{HTML}{173B57}
\definecolor{agiGold}{HTML}{B7791F}
\definecolor{agiLight}{HTML}{EFF6FA}
\definecolor{agiGreen}{HTML}{184D3B}
\hypersetup{colorlinks=true,linkcolor=agiBlue,urlcolor=agiBlue}
\setlength{\parindent}{0pt}
\setlength{\parskip}{0.55em}
\setlength{\headheight}{14pt}
\setlist{nosep,leftmargin=*}
\pagestyle{fancy}
\fancyhf{}
\fancyhead[L]{\small Part of the \href{https://github.com/mmikhasenko/GIModel.jl}{Agentic Gottried--Isgur} project}
\fancyhead[R]{\SheetCode}
\fancyfoot[C]{\thepage}
\newcommand{\dd}{\,\mathrm d}
\newcommand{\ket}[1]{\lvert #1\rangle}
\newcommand{\bra}[1]{\langle #1\rvert}
\newcommand{\braket}[2]{\langle #1\mid #2\rangle}
\newcommand{\expect}[1]{\left\langle #1\right\rangle}
\newcommand{\op}[1]{\widehat{#1}}
\newcommand{\GeV}{\,\mathrm{GeV}}
\newcommand{\R}{\mathbb{R}}
\newtcolorbox{goals}{colback=agiLight,colframe=agiBlue,title=Learning outcomes}
\newtcolorbox{reading}{colback=white,colframe=agiGold,title=Book reference,
  after skip=5mm}
\newtcolorbox{paperbridge}{colback=agiGreen!7,colframe=agiGreen,
  colbacktitle=agiGreen,coltitle=white,title=Connection to the GI paper}
\newtcolorbox{problem}[1]{colback=white,colframe=agiBlue,title={#1},
  before skip=3.5mm,after skip=1mm}
\newtcolorbox{solution}{colback=black!2,colframe=black!45,title=Solution}
\newtcolorbox{quiz}{colback=white,colframe=agiGold,title=Post-solution concept check}
\newtcolorbox{checkpoint}{colback=white,colframe=agiGold,title=Interpretation checkpoint}
\newcommand{\sheettitle}[3]{%
  \begin{center}
  {\Large\bfseries #1}\par
  \vspace{0.2em}{\color{agiBlue}#2}\par
  \vspace{0.4em}\small #3
  \end{center}\vspace{0.5em}}
\newcommand{\difficulty}[1]{\textbf{Difficulty:} #1}
\newcommand{\timebox}[1]{\textbf{Suggested time:} #1}
\newcommand{\answer}{\par\smallskip\color{black!50}\textbf{Answer:} }

\begin{document}
\sheettitle{Sheet 3: From a Hamiltonian to a matrix}{Variational mechanics, oscillator bases, and $\sqrt{p^2+m^2}$}{Prerequisite: Sheets 1--2. \difficulty{★/★★★} \quad \timebox{4--5 hours}}

\begin{goals}
Derive Rayleigh--Ritz; understand the harmonic-oscillator (HO) basis as a
representation rather than a physical assumption; prove variational monotonicity;
analyze the semirelativistic kinetic energy and the role of a basis scale $\beta$;
construct the oscillator radial functions explicitly and obtain their $r^2$ and
$p^2$ matrix elements by two independent analytic routes.
\end{goals}
\begin{reading}
Selected reading in Landau--Lifshitz, \emph{Quantum Mechanics: Non-Relativistic
Theory}: Ch. III, ``Schr\"odinger's Equation,'' for the variational principle,
and Ch. VI, ``Perturbation Theory.'' These chapters are optional background;
the sheet develops the required argument directly.
\end{reading}


\begin{problem}{Problem 1 ★ — Derive Rayleigh--Ritz}
For $\ket\psi=\sum_{i=1}^N c_i\ket{i}$ in an orthonormal basis, impose
$\braket\psi\psi=1$ and derive the stationary equation for
$\expect H$.
\end{problem}
\begin{solution}
Stationarity of $c^\dagger Hc-E(c^\dagger c-1)$ with respect to $c_i^*$ gives
$Hc=Ec$. The matrix $H_{ij}=\bra iH\ket j$ is the Hamiltonian restricted, or
projected, to the span of the chosen $N$ basis functions. Diagonalizing it
therefore solves the constrained variational problem in that finite subspace.
\end{solution}
\begin{quiz}
Does matrix diagonalization approximate the algebra or the Hilbert space?
\answer the finite subspace approximates the Hilbert space; diagonalization
solves that projected problem exactly up to arithmetic error.
\end{quiz}

\begin{problem}{Problem 2 ★★ — Prove monotone upper bounds}
Let $V_N\subset V_{N+1}$. Prove that the lowest variational energy (also called
the lowest Ritz value) obeys
$E_0^{(N+1)}\le E_0^{(N)}$ and $E_0^{(N)}\ge E_0$.
\end{problem}
\begin{solution}
The minimum of the Rayleigh quotient over a larger set cannot increase, giving
the first inequality. The exact ground energy is the minimum over the full
domain, which contains $V_N$, giving the second. Corresponding upper bounds for
excited states follow when the minimization is restricted to states orthogonal
to the already determined lower states.
\end{solution}
\begin{quiz}
Why are two successive nearly equal energies evidence, but not a proof, of
convergence? \answer the result may be insensitive to the last change in $N$
while still depending on an untested basis scale $\beta$, or an excited level
may converge more slowly. More than one independent numerical control must be
varied.
\end{quiz}

\begin{problem}{Problem 3 ★ — Expand the relativistic kinetic energy}
Expand $T(p)=\sqrt{p^2+m_1^2}+\sqrt{p^2+m_2^2}$ through $p^4$. Identify the
nonrelativistic reduced-mass term and state the expansion parameter.
\end{problem}
\begin{solution}
\[
T=m_1+m_2+\frac{p^2}{2}\left(\frac1{m_1}+\frac1{m_2}\right)
-\frac{p^4}{8}\left(\frac1{m_1^3}+\frac1{m_2^3}\right)+O(p^6/m^5).
\]
Since $1/\mu=1/m_1+1/m_2$, the quadratic term is $p^2/(2\mu)$. The expansion
requires typical $p^2/m_i^2\ll1$, which is least reliable for light quarks.
\end{solution}
\begin{quiz}
Why not truncate at $p^4$ in a model intended to span pion to upsilon? \answer
there is no uniformly small light-quark expansion parameter.
\end{quiz}

\begin{problem}{Problem 4 ★★ — Use a Gaussian trial state}
For the normalized three-dimensional Gaussian
$\psi_\beta=(\beta^3/\pi^{3/2})^{1/2}e^{-\beta^2r^2/2}$, compute
$\expect{r}$, $\expect{1/r}$, and $\expect{p^2}$.
\end{problem}
\begin{solution}
Radial Gaussian integrals give
\[
\expect r=\frac{2}{\sqrt\pi\,\beta},\qquad
\expect{1/r}=\frac{2\beta}{\sqrt\pi},\qquad
\expect{p^2}=\frac{3\beta^2}{2}.
\]
For $V=br-\kappa/r$, increasing $\beta$ lowers the Coulomb contribution but
raises the kinetic energy and reduces the confinement energy.
\end{solution}
\begin{quiz}
What physical length does $\beta^{-1}$ encode? \answer the trial state's radial
size; optimizing $\beta$ adapts the representation to the bound state.
\end{quiz}


\begin{problem}{Problem 5 ★★ — Build and normalize the oscillator basis}
The reduced radial oscillator functions at fixed $L$ are
\[
u_{nL}(r)=N_{nL}\,(\beta r)^{L+1}e^{-(\beta r)^2/2}\,
L_n^{L+1/2}\!\left((\beta r)^2\right),\qquad n=0,1,2,\dots
\]
where $L_n^\alpha$ is an associated Laguerre polynomial. Using only
\[
\int_0^\infty x^\alpha e^{-x}L_n^\alpha(x)L_m^\alpha(x)\dd x
=\frac{\Gamma(n+\alpha+1)}{n!}\,\delta_{nm},
\]
fix $N_{nL}$ from $\int_0^\infty u_{nL}u_{mL}\dd r=\delta_{nm}$, and write
$u_{0L}$ and $u_{1L}$ out explicitly. Which two conventions does normalization
leave unfixed?
\end{problem}
\begin{solution}
Substitute $x=(\beta r)^2$, so that $\dd r=\dd x/(2\beta\sqrt x)$ and
$(\beta r)^{2L+2}=x^{L+1}$. With $\alpha=L+\tfrac12$ the measure combines into
$x^{\alpha}\dd x/(2\beta)$ and
\[
\int_0^\infty u_{nL}u_{mL}\dd r
=\frac{N_{nL}N_{mL}}{2\beta}\int_0^\infty x^{\alpha}e^{-x}
L_n^\alpha L_m^\alpha\dd x
=\frac{N_{nL}^2}{2\beta}\frac{\Gamma(n+L+\tfrac32)}{n!}\,\delta_{nm},
\]
so orthogonality is inherited from the Laguerre family and
\[
N_{nL}=\sqrt{\frac{2\beta\,n!}{\Gamma(n+L+\tfrac32)}}.
\]
Since $L_0^\alpha=1$ and $L_1^\alpha(x)=1+\alpha-x$,
\[
u_{0L}=\sqrt{\frac{2\beta}{\Gamma(L+\tfrac32)}}\,(\beta r)^{L+1}e^{-x/2},
\qquad
u_{1L}=\sqrt{\frac{2\beta}{\Gamma(L+\tfrac52)}}\,(\beta r)^{L+1}e^{-x/2}
\left(L+\tfrac32-x\right).
\]
Normalization fixes $|N_{nL}|$ but not its sign, and it does not record whether
the polynomial argument is $\beta r$ or $(\beta r)^2$. Both conventions must be
stated before any analytic matrix element is written down.
\end{solution}
\begin{quiz}
The textbook oscillator carries a mass $\mu$ and a frequency $\omega$, yet
neither appears above. Where did they go? \answer into the single combination
$\beta^2=\mu\omega$. Here it is a variational parameter of the representation,
not a physical frequency of the meson, which is why Problem 7 below may scan it
freely.
\end{quiz}

\begin{problem}{Problem 6 ★★★ — Two independent routes to $r^2$ and $p^2$}
(a) Using the three-term recurrence
\[
xL_m^\alpha(x)=(2m+\alpha+1)L_m^\alpha-(m+1)L_{m+1}^\alpha-(m+\alpha)L_{m-1}^\alpha
\]
together with the orthogonality relation of Problem 5, compute every matrix element
$\bra{n}r^2\ket{m}$ in that basis.
(b) Now obtain $\bra{n}p^2\ket{m}$ with no integral at all, using only the fact
that $H_{\rm osc}=p^2/(2\mu)+\tfrac12\mu\omega^2r^2$ is diagonal in this basis
with eigenvalue $(2n+L+\tfrac32)\omega$.
(c) How many independent numbers does each $N\times N$ matrix contain?
\end{problem}
\begin{solution}
(a) Write $h_n=\Gamma(n+\alpha+1)/n!$, so $N_{nL}^2=2\beta/h_n$, and use
$r^2=x/\beta^2$. Applying the recurrence to $xL_m^\alpha$ and projecting leaves
only $n=m$ and $n=m\pm1$. The diagonal term carries $(2m+\alpha+1)h_m$, and the
raising term carries $-(m+1)h_{m+1}$; with
$h_{m+1}/h_m=(m+\alpha+1)/(m+1)$ the normalizations collapse to
\[
\bra{n}r^2\ket{n}=\frac{2n+L+\tfrac32}{\beta^2},\qquad
\bra{n+1}r^2\ket{n}=-\frac{\sqrt{(n+1)(n+L+\tfrac32)}}{\beta^2},
\]
every other element vanishing. The operator $r^2$ is symmetric tridiagonal.

(b) With $\beta^2=\mu\omega$ the oscillator relation rearranges to
$p^2=2\mu H_{\rm osc}-\beta^4r^2$. The first term is diagonal with
$2\beta^2(2n+L+\tfrac32)$, and subtracting $\beta^4$ times the result of (a)
gives
\[
\bra{n}p^2\ket{n}=\beta^2\left(2n+L+\tfrac32\right),\qquad
\bra{n+1}p^2\ket{n}=+\beta^2\sqrt{(n+1)(n+L+\tfrac32)} .
\]
Same magnitudes as $r^2$ up to $\beta^{\pm2}$, opposite off-diagonal sign, and
not one integral was evaluated.

(c) Each is symmetric tridiagonal: $N$ diagonal plus $N-1$ off-diagonal
entries, so $2N-1$ numbers rather than $N^2$ integrals. This linear count, not
raw computing power, is what made a large-basis diagonalization affordable long
before it was cheap to integrate anything numerically.
\end{solution}
\begin{quiz}
The off-diagonal elements of $r^2$ are negative and those of $p^2$ positive,
while the diagonals agree up to $\beta^{\pm2}$. Is one of the two a sign error?
\answer no. $p^2=2\mu H-\beta^4r^2$ subtracts the $r^2$ matrix from a diagonal
one, which flips precisely the off-diagonal sign. Replacing $u_{nL}\to-u_{nL}$
for some $n$ flips the corresponding off-diagonals of \emph{both} matrices
together, so a convention that appears to flip only one of them means an
analytic formula is being combined with a differently phased basis. Compare
\texttt{ho\_r2\_matrix} and \texttt{ho\_p2\_matrix} in
\texttt{src/harmonic\_oscillator\_basis.jl}.
\end{quiz}

\begin{problem}{Problem 7 ★★★ — Explain the order of the calculation in the paper}
Suppose $H(\beta,N)$ is the GI matrix for one fixed choice of $(L,S,J)$ in an $N$-function HO
basis, built from the functions of Problem 5. Design a pen-and-paper convergence logic using $N$, $N+\Delta N$, and
$\beta$. Why must mixing between different $(L,S,J)$ choices be applied only
after these separately calculated states converge?
\end{problem}
\begin{solution}
For each $N$, minimize the target variational energy over a chosen interval of
$\beta$; enlarge
$N$ and repeat until two successive changes are below tolerance and the optimum
does not occur at an endpoint of that interval. Retain enough converged levels for every
later mixing block. Mixing before convergence entangles truncation error with a
physical off-diagonal coupling: changing $N$ then changes both the basis states
and the matrix being interpreted.
\end{solution}
\begin{quiz}
If HO energies agree with those from a finite-difference method, which represents
$u(r)$ by values at selected radii, but their stored arrays look different, is
that a contradiction? \answer no. Coefficients in basis functions and values at
radii are different representations; norms, matrix elements, and eigenvalues
are the quantities to compare.
\end{quiz}

\begin{paperbridge}
GI Eq. (14) collects the relativized kinetic, confinement, hyperfine, and
spin--orbit operators. The paragraph following it states that this Hamiltonian
is directly diagonalized in large harmonic-oscillator bases, separately in
$\ket{jm;ls}$ sectors, and that the basis is expanded until convergence.
Eq. (A17) is where Problems 5--6 are actually used: the momentum-dependent
factors act on exactly known $p^2$ matrix elements while the radial potentials
are evaluated in the same oscillator basis, with no spatial mesh anywhere. Those
two problems are the analytic content of
\texttt{src/harmonic\_oscillator\_basis.jl}; Sheet 5 supplies the potentials
that go between them.
\end{paperbridge}
\end{document}
